\section{From the Exact Quantity to Observable Diagnostics}\label{sec:diagnostics}

Computing $\mathsf{CD}_\tau(\Pi)$ requires the joint law of the microstate and the packaged future. In applications, one usually has less: an induced macro kernel, a packaged time series, or a predictor of limited size. This section connects the deficit to three such handles, summarized in \Cref{tab:diagnostics}. Each comes with a precise statement of what it does and does not certify.

\begin{table}[tbp]
\centering
\caption{The exact quantity and the three handles used to approach it. Each row names the result in this section that connects it to the deficit.}
\label{tab:diagnostics}
\small
\setlength{\tabcolsep}{5pt}
\renewcommand{\arraystretch}{1.15}
\begin{tabularx}{\linewidth}{@{}>{\RaggedRight\arraybackslash}p{0.17\linewidth} >{\RaggedRight\arraybackslash}X >{\RaggedRight\arraybackslash}p{0.30\linewidth}@{}}
\toprule
Quantity & What it measures & Relation to $\mathsf{CD}_\tau(\Pi)$ \\
\midrule
$\mathsf{CD}_\tau(\Pi)$ & information about the next package hidden inside the current one & target (\Cref{def:closure-deficit}) \\
$\mathrm{RM}_\tau^{\mathrm{stat}}(\Pi)$ & average $\ell_1$ distance between microstate forecasts and package forecasts & $\mathsf{CD}\ge\tfrac12(\mathrm{RM}^{\mathrm{stat}})^2$ (\Cref{prop:rm-lower-bound}) \\
$H(Y_{t+\tau}\mid Y_t)-H(Y_{t+\tau}\mid S_L)$ & predictive gain from $L$ staged packages of memory & between $0$ and $\mathsf{CD}$ for Markov substrates (\Cref{prop:memory}) \\
$\mathsf{Rand}_B(\Pi,\tau)$ & best expected log loss within a budgeted predictor class & equals $H(Y_{t+\tau}\mid S^{(B)})$ for a full class (\Cref{prop:budget}) \\
\bottomrule
\end{tabularx}
\end{table}

\subsection{Route mismatch}

Coarse-graining pipelines often summarize a packaging by an induced macro kernel. One chooses a \emph{lift}, a probability distribution over the microstates of each package. One then runs the micro dynamics for $\tau$ steps from that distribution and records the resulting distribution of packages \citep{tsiokos2026pica}. Call the resulting row $\widehat p_y^{(\tau,\mathrm{lift})}$. Its \emph{route mismatch} measures how far the true microstate forecasts are from it:
\begin{equation}
\mathrm{RM}_\tau^{\mathrm{lift}}(\Pi)
:=
\sum_{x \in \mathcal{X}}
\pi(x)\,
\bigl\|
p_x^{(\tau)} - \widehat p_{\Pi(x)}^{(\tau,\mathrm{lift})}
\bigr\|_1 .
\label{eq:route-mismatch}
\end{equation}
Two lifts matter here. The \emph{stationary lift} weights the microstates of package $y$ by $\pi(x\mid y)$; by \eqref{eq:fiber-average}, its row is exactly the package forecast $\bar p_y^{(\tau)}$, and we write $\mathrm{RM}_\tau^{\mathrm{stat}}$. The \emph{uniform lift} weights the microstates of each package equally, giving $\mathrm{RM}_\tau^{\mathrm{unif}}$. Both versions average over microstates with the stationary weights $\pi(x)$. Only the macro rows differ.

\begin{proposition}[Route-mismatch bound]\label{prop:rm-lower-bound}
\begin{equation}
\mathsf{CD}_\tau(\Pi)
\ge
\frac{1}{2}\bigl(\mathrm{RM}_\tau^{\mathrm{stat}}(\Pi)\bigr)^2.
\label{eq:rm-lower-bound}
\end{equation}
\end{proposition}

\begin{proof}
Pinsker's inequality in nats states that $D_{\mathrm{KL}}(p\|q)\ge\tfrac12\|p-q\|_1^2$ for probability vectors $p,q$ \citep{cover2006elementsofinformationtheory}. Apply it to each term of \eqref{eq:cd-kl}:
\[
\mathsf{CD}_\tau(\Pi)
\ge
\frac12\sum_{x:\,\pi(x)>0}\pi(x)\,\bigl\|p_x^{(\tau)}-\bar p_{\Pi(x)}^{(\tau)}\bigr\|_1^2
\ge
\frac12\Bigl(\sum_{x}\pi(x)\,\bigl\|p_x^{(\tau)}-\bar p_{\Pi(x)}^{(\tau)}\bigr\|_1\Bigr)^2 .
\]
The second step is Jensen's inequality for the convex map $u\mapsto u^2$ under the probability weights $\pi$; unoccupied states contribute nothing to either side. The last sum is $\mathrm{RM}_\tau^{\mathrm{stat}}(\Pi)$.
\end{proof}

In total-variation units, $\|p-q\|_{\mathrm{TV}}=\tfrac12\|p-q\|_1$, the same bound reads $\mathsf{CD}\ge 2\,(\mathrm{RM}^{\mathrm{stat}}_{\mathrm{TV}})^2$. The bound is one-sided: a small deficit forces a small stationary mismatch, but Pinsker's inequality alone gives no upper bound on the deficit in terms of the mismatch. (Entropy-continuity arguments do give such a bound, with a constant that depends on the size of the package alphabet.)

\begin{remark}[The uniform lift is not covered]\label{rem:uniform-lift}
The uniform lift is common in practice because it does not require the stationary law to build the macro rows \citep{tsiokos2026pica}. It does not inherit \Cref{prop:rm-lower-bound}. In the three-state chain of \Cref{ex:uniform-lift}, with a rarely visited microstate, the deficit is about $0.0040$ nats. The uniform-lift mismatch is about $0.500$, so the would-be bound $\tfrac12(\mathrm{RM}^{\mathrm{unif}})^2\approx 0.125$ fails by a factor of thirty. The stationary mismatch, about $0.0020$, satisfies the true bound. The uniform-lift mismatch is still a reasonable \emph{proxy}, and it tracks the deficit closely in our benchmark (\Cref{sec:results-markov}). That agreement is an empirical observation, not a theorem.
\end{remark}

\subsection{Packaged memory}\label{sec:memory}

Often the microstate is never observed, and only the packaged sequence is available. One can still ask whether remembering earlier packages improves prediction. Write the staged packaged history of length $L\ge1$ as
\[
S_L := \bigl(Y_{t-(L-1)\tau},\dots,Y_{t-\tau},Y_t\bigr),
\]
so that $S_1=Y_t$. The next result shows that, for a Markov substrate, packaged memory can only recover part of the deficit.

\begin{proposition}[Memory is sandwiched by the deficit]\label{prop:memory}
Let $(X_t)$ be a stationary Markov chain. For every $L\ge1$,
\begin{equation}
H(Y_{t+\tau}\mid X_t)
\;\le\;
H(Y_{t+\tau}\mid S_{L+1})
\;\le\;
H(Y_{t+\tau}\mid S_L)
\;\le\;
H(Y_{t+\tau}\mid Y_t).
\label{eq:memory-sandwich}
\end{equation}
Consequently,
\begin{equation}
0
\;\le\;
H(Y_{t+\tau}\mid Y_t)-H(Y_{t+\tau}\mid S_L)
\;\le\;
\mathsf{CD}_\tau(\Pi).
\label{eq:memory-gain}
\end{equation}
\end{proposition}

\begin{proof}
The history $S_L$ is a sub-vector of $S_{L+1}$, and $Y_t$ is a component of $S_L$. Conditioning on more variables cannot increase entropy, which gives the middle and right inequalities. For the left inequality, $S_{L+1}$ is a function of $(X_{t-L\tau},\dots,X_t)$. By the Markov property, $Y_{t+\tau}$ is conditionally independent of $(X_s)_{s<t}$ given $X_t$. Hence $H(Y_{t+\tau}\mid X_t,S_{L+1})=H(Y_{t+\tau}\mid X_t)$, and
$H(Y_{t+\tau}\mid S_{L+1})\ge H(Y_{t+\tau}\mid X_t,S_{L+1})=H(Y_{t+\tau}\mid X_t)$.
The bound \eqref{eq:memory-gain} follows by subtracting from $H(Y_{t+\tau}\mid Y_t)$ and using \Cref{prop:cd-decomposition}.
\end{proof}

Two consequences are worth stating. First, the deficit is a hard ceiling on what any finite window of packaged memory can buy back. The rest is information that the packaged past does not carry. Second, the one-step version gives a clean reading of the predictive-gap diagnostic used in earlier Six Birds work \citep{tsiokos2026notch}.

\begin{corollary}[Population predictive gap]\label{cor:pred-gap}
For a stationary Markov substrate, the population gap between first- and second-order packaged prediction is
\[
\Delta^{(2,1)}
:=
H(Y_{t+\tau}\mid Y_t)-H(Y_{t+\tau}\mid Y_{t-\tau},Y_t)
=
I(Y_{t-\tau};Y_{t+\tau}\mid Y_t),
\]
and it satisfies $0\le\Delta^{(2,1)}\le\mathsf{CD}_\tau(\Pi)$.
\end{corollary}

\begin{proof}
The equality is the definition of conditional mutual information. The bounds are \eqref{eq:memory-gain} with $L=2$.
\end{proof}

A positive population gap therefore certifies a positive deficit. A zero gap does not certify closure, because the hidden distinction may leave no trace in the packaged past. In practice, one computes the difference of held-out average log losses of \emph{fitted} order-one and order-two predictors. That difference estimates the difference of the fitted predictors' risks, which equals $\Delta^{(2,1)}$ plus the difference of their estimation errors (each fitted predictor's average KL divergence from the true conditional law). It approximates $\Delta^{(2,1)}$ when both fits are accurate, and on finite data it can even be negative.

\subsection{Budgeted randomness}

The memory picture suggests pricing randomness by budget. Let $\mathcal{Q}_B$ be the class of predictors available under budget $B$. Each predictor in the class maps a summary $S^{(B)}_t$ of the observable past to a distribution over the next package.

\begin{definition}[Budgeted randomness]\label{def:budgeted-rand}
\begin{equation}
\mathsf{Rand}_B(\Pi,\tau)
:=
\inf_{Q \in \mathcal{Q}_B}
\mathbb{E}\!\left[
-\log Q\bigl(Y_{t+\tau}\mid S_t^{(B)}\bigr)
\right].
\label{eq:budgeted-rand}
\end{equation}
\end{definition}

\begin{proposition}[Properties of budgeted randomness]\label{prop:budget}
\begin{enumerate}
\item[(i)] For every $Q$, the expected log loss equals $H(Y_{t+\tau}\mid S^{(B)}_t)$ plus the average KL divergence from the true conditional law to $Q$. Hence $\mathsf{Rand}_B\ge H(Y_{t+\tau}\mid S^{(B)}_t)$, with equality when $\mathcal{Q}_B$ contains every conditional distribution given $S^{(B)}_t$.
\item[(ii)] Suppose budgets are nested: whenever $B\le B'$, the summary $S^{(B)}_t$ is a function of $S^{(B')}_t$, and every predictor in $\mathcal{Q}_B$, composed with that function, belongs to $\mathcal{Q}_{B'}$. Then $\mathsf{Rand}_{B'}\le\mathsf{Rand}_B$.
\item[(iii)] For a stationary Markov substrate, with memory budget $B=L$, summary $S^{(L)}_t=S_L$, and full predictor classes, $\mathsf{Rand}_L=H(Y_{t+\tau}\mid S_L)$. This value is nonincreasing in $L$ and lies between $H(Y_{t+\tau}\mid X_t)$ and $H(Y_{t+\tau}\mid Y_t)$.
\end{enumerate}
\end{proposition}

\begin{proof}
(i) is the standard decomposition of log loss into conditional entropy and divergence; the divergence vanishes for the true conditional law. (ii) holds because the infimum for $B'$ is taken over a set that contains a copy of every predictor available at $B$, with the same loss. (iii) combines (i) with \Cref{prop:memory}.
\end{proof}

Monotonicity in (ii) is a property of nested classes, not of budgets as such. If enlarging the budget removes options, for example by forcing a different summary, the curve can rise. Moreover, $\mathsf{Rand}_B$ is a \emph{population} quantity. The held-out loss of a fitted model estimates that model's risk, which is at least $\mathsf{Rand}_B$ and approaches it only when the fit is accurate; on a finite test set the estimate itself fluctuates. Choosing the best of several fitted models by their test loss produces a curve that decreases by construction. \Cref{sec:results-budget} therefore reports exact population entropies alongside models selected on separate validation data.

\paragraph{What the diagnostics do not measure.} A positive deficit does not imply that a new theory must be born, and it does not certify physical irreversibility. Arrow-of-time diagnostics compare forward and time-reversed path measures and behave differently under coarse-graining \citep{esposito2012stochasticthermodynamicsundercoarsegraining,tsiokos2026sixbirds}. A hidden protocol phase or an omitted variable can make a layer look random without any underlying nonequilibrium drive. Route mismatch and the predictive gap are views of the deficit, not substitutes for it.
